\documentclass[11pt]{article}
\usepackage[margin=1in]{geometry}
\usepackage{amsmath,amssymb}
\usepackage[T1]{fontenc}
\usepackage{lmodern}
\usepackage{hyperref}

\title{Literature answer to a uniform-boundedness problem in limit operators}
\author{Run fa\_banach\_001}
\date{2026-06-28}

\begin{document}
\maketitle
\sloppy

\paragraph{Status.}
This is a \textbf{literature already answered} packet for Open Problem 8 of
Chandler-Wilde--Lindner, not a new proof packet.

\paragraph{Source problem.}
Chandler-Wilde--Lindner, \emph{Limit Operators, Collective Compactness, and the
Spectral Theory of Infinite Matrices} (arXiv:1005.0166), end with a list of
open problems.  Open Problem 8 asks whether the uniform boundedness condition
in Theorem \texttt{prop\_rich\_invinf}(iii) is redundant for
\(p\in(1,\infty)\).  In parsed source lines 7306--7312, the paper says that
the redundancy was known for \(p\in\{0,1,\infty\}\), while for the remaining
cases \(p\in(1,\infty)\) the question remained open.

\paragraph{Supporting answer.}
Lindner--Seidel, \emph{An Affirmative Answer to a Core Issue on Limit Operators}
(arXiv:1401.1300), directly addresses the same issue.  Its abstract says that
for rich band-dominated operators the known Fredholm criterion required all
limit operators to be invertible with uniformly bounded inverses, and that the
paper shows the uniform-boundedness condition is redundant.

The introduction identifies the question as the ``big question'': whether the
operator spectrum of a rich operator is automatically uniformly invertible if
it is elementwise invertible.  Lines 265--267 of the parsed source state that
the aim is to solve this problem in general by showing that the condition is
redundant.  The precise statements are Corollary 3.7 and Theorem 3.8.  Corollary
3.7 proves that if \(A\) is rich and band-dominated and all
\(B\in\sigma_{\mathrm{op}}(A)\) are invertible, then the norms
\(\|B^{-1}\|\) are uniformly bounded.  Theorem 3.8 then gives the simplified
criterion:
\[
  A\in\mathcal A_p^\$, \quad
  A \text{ is } \mathcal P\text{-Fredholm}
  \iff
  \text{all limit operators of } A \text{ are invertible},
\]
for \(A\) acting on \(\ell^p(\mathbb Z^N,X)\) with
\(p\in\{0\}\cup[1,\infty]\).

\paragraph{Identification and scope.}
This answers the source's Open Problem 8 affirmatively, and in particular
covers the previously open range \(1<p<\infty\).  The answer is literature,
not new work by this run.  The supporting theorem keeps the rich and
band-dominated hypotheses; Section 4 of the supporting paper gives examples
showing these hypotheses cannot simply be dropped.

\paragraph{Search evidence.}
The cheap run indexes were searched for arXiv:1005.0166, arXiv:1401.1300,
limit operators, collective compactness, Fredholmness, essential spectrum,
and the uniform boundedness condition.  No existing run packet recorded this
exact answer.  External arXiv search for phrases around ``uniform boundedness
condition'', ``limit operators'', and ``redundant'' located arXiv:1401.1300 as
the decisive later source.

\begin{thebibliography}{9}
\bibitem{ChandlerWildeLindner2010}
S. N. Chandler-Wilde and M. Lindner,
\emph{Limit Operators, Collective Compactness, and the Spectral Theory of
Infinite Matrices},
arXiv:1005.0166.

\bibitem{LindnerSeidel2014}
M. Lindner and M. Seidel,
\emph{An Affirmative Answer to a Core Issue on Limit Operators},
J. Funct. Anal. 267(3) (2014), 901--917; arXiv:1401.1300.
\end{thebibliography}

\end{document}
